% !TeX root = ../../negocio-tecnologico-exemplo-undf.tex
\undfsetcontent{resumo}{%
Este relatório apresenta uma análise didática de modelo de negócio para um
serviço digital de apoio à organização de trabalhos acadêmicos. O objetivo
foi construir uma proposta documentada, identificar hipóteses de funcionamento
e examinar sua consistência em um cenário financeiro simplificado. O método
compreendeu consulta documental sobre a ferramenta LaTeX, organização de nove
componentes de negócio e cálculo de três cenários com premissas inteiramente
hipotéticas. O serviço foi delimitado ao suporte de composição documental,
preservando a autoria intelectual do estudante. A investigação identificou
uma base técnica disponível, mas não forneceu evidências de disposição a
pagar, demanda ou capacidade operacional. Com preço hipotético de noventa
reais, custo variável unitário de trinta reais e custo fixo de seiscentos
reais por período, o equilíbrio aritmético ocorreu em dez atendimentos.
Esse resultado não comprova viabilidade comercial e não incorpora todos os
custos de uma operação real. Não houve entrevistas, vendas, desenvolvimento
de plataforma ou implantação de empresa. Concluiu-se que a proposta possui
hipóteses documentadas, mas carece de evidências para recomendar sua implantação.
O exemplo demonstra a entrega acadêmica final de uma análise, incluindo
resultados inconclusivos e recomendações decorrentes de seus limites.}
\undfsetcontent{palavraschave}{modelo de negócio; base tecnológica; análise de viabilidade}
\undfsetcontent{abstract}{%
This report presents an educational analysis of a business model for a digital
service supporting the organization of academic documents. The objective was
to produce a documented proposal, identify operating assumptions and examine
its consistency within a simplified financial scenario. The method included
documentary research on LaTeX, the organization of nine business components
and calculations for three scenarios using entirely hypothetical assumptions.
The service was limited to document composition support, preserving the
student's intellectual authorship. The investigation identified an available
technical foundation but provided no evidence of willingness to pay, demand
or operational capacity. With a hypothetical price of ninety reais, a variable
unit cost of thirty reais and fixed costs of six hundred reais per period,
arithmetic break-even occurred at ten service units. This result does not
demonstrate commercial feasibility and does not include every cost of a real
operation. No interviews, sales, platform development or company deployment
took place. The analysis concluded that the proposal has documented hypotheses
but lacks evidence to recommend implementation. The example demonstrates a
completed academic analysis, including inconclusive findings and recommendations
derived from its limitations.}
\undfsetcontent{keywords}{business model; technology-based business; feasibility analysis}
